% ======================================================================
% 10-trace-to-zeta.tex (FULL REWRITE, Annals+++ / DIAMOND)
% Chapter 10. Trace-to-zeta bridges: determinants, explicit formulae,
%            and positivity gates
%
% Design principle (Annals+++):
%   - Every bridge is an OBJECT + a STABILITY STATEMENT + a LEDGER BOUND.
%   - Identity / continuous-spectrum terms are isolated explicitly (no hidden swaps).
%   - Positivity is expressed as a quadratic form with an explicit correction + ledger remainder.
%   - All auxiliary dependence is confined to: (i) explicit identity correction, (ii) Chapter-8 ledger.
% ======================================================================

\chapter{Trace-to-zeta bridges: determinants, explicit formulae, and positivity criteria}
\label{chap:10-trace-to-zeta}

\epigraph{A trace formula becomes a zeta function only after normalization has turned ``test dependence'' into an object and an error ledger.}{}

% ======================================================================
\section{Overview and bridge taxonomy}
\label{sec:10.overview}
% ======================================================================

This chapter records, in a form compatible with the normalized trace framework of
Chapter~\ref{chap:08-normalization}, three standard analytic interfaces between trace identities and zeta-type objects:

\begin{enumerate}
\item \textbf{Determinant bridge:} from stabilized trace data to regularized determinants and logarithmic derivatives;
\item \textbf{Explicit-formula bridge:} from the normalized trace identity to a canonical spectral--geometric dictionary;
\item \textbf{Positivity bridge:} from convolution squares to quadratic forms whose nonnegativity constrains spectral parameters.
\end{enumerate}

The emphasis is on \emph{portable} formulations: every statement is invariant under admissible auxiliary modifications
up to an explicit identity-sector correction and a remainder controlled by the consolidated ledger of Chapter~8.

\medskip
\noindent\textbf{Standing lock.}
Throughout, $\mathcal T_\Gamma[\cdot]$ denotes the normalized trace functional
\[
\mathcal T_\Gamma[h]=\mathcal L_\Gamma[h-h(0)\chi]
\]
for a fixed admissible cutoff $\chi$ as in Chapter~\ref{chap:08-normalization}; any change of admissible construction is
controlled by Theorem~\ref{thm:08.normalization-stability}.

% ======================================================================
\section{Regularized determinants from stabilized trace input}
\label{sec:10.det}
% ======================================================================

\subsection{Spectral zeta functions and zeta-regularized determinants}
\label{subsec:10.det.zeta}

Let $A$ be a self-adjoint operator on a separable Hilbert space $H$, bounded below, with compact resolvent.
Let $\{\lambda_j\}_{j\ge1}$ be its nonzero eigenvalues, repeated with multiplicity.
For $\Re(s)$ sufficiently large define
\begin{equation}\label{eq:10.zetaA.def}
\zeta_A(s):=\sum_{j\ge1}\lambda_j^{-s}.
\end{equation}
Assume $\zeta_A(s)$ admits a meromorphic continuation to a neighborhood of $s=0$ that is holomorphic at $s=0$.
Then the zeta-regularized determinant is
\begin{equation}\label{eq:10.detA.def}
\det\nolimits_{\zeta}(A):=\exp\!\bigl(-\zeta_A'(0)\bigr).
\end{equation}

A standard construction uses heat traces: if $e^{-tA}$ is trace class for all $t>0$ and $\Tr(e^{-tA})$ admits a small-time
asymptotic expansion, then for $\Re(s)$ large,
\begin{equation}\label{eq:10.mellin}
\zeta_A(s)=\frac{1}{\Gamma(s)}\int_0^\infty t^{s-1}\bigl(\Tr(e^{-tA})-\Pi_0\bigr)\,dt,
\end{equation}
where $\Pi_0$ is the nullspace contribution. Analytic continuation of \eqref{eq:10.mellin} yields the continuation of $\zeta_A$.

\subsection{Trace formula input as a normalized object}
\label{subsec:10.det.trace-input}

In the geometric settings of interest the trace formula provides distributional identities of the schematic form
\begin{equation}\label{eq:10.trace.distributional}
\Tr^{\mathrm{reg}}\!\bigl(\varphi(A)\bigr)
=
\mathcal I[\varphi]+\mathcal G[\varphi]+\mathcal C[\varphi],
\end{equation}
where $\mathcal I$ is the identity term, $\mathcal G$ the geometric (orbital) term, and $\mathcal C$ the continuous-spectrum
(scattering) term (vanishing in compact cases).

Normalization (Chapter~\ref{chap:08-normalization}) replaces \eqref{eq:10.trace.distributional} by a canonical trace input
\begin{equation}\label{eq:10.trace.normalized.input}
\mathcal T_\Gamma[h]=\mathcal G_\Gamma[h]+\mathcal S_\Gamma[h]+\Err_\Gamma[h],
\end{equation}
where $\mathcal S_\Gamma$ denotes the remaining spectral contribution after identity subtraction (including continuous spectrum
in the fixed regularization), and $\Err_\Gamma[h]$ is controlled by the consolidated ledger $\mathcal E_\Gamma[h]$.

\medskip
\noindent\textbf{Determinant principle (portable form).}
Determinant-level objects may be reconstructed from trace data only through inputs of the type \eqref{eq:10.trace.normalized.input},
so that all admissible auxiliary dependence is confined to (i) an explicit identity correction and (ii) ledger-bounded errors.

% ======================================================================
\section{Smoothed oscillatory traces as spectral probes}
\label{sec:10.osc}
% ======================================================================

\subsection{A standard localized test family}
\label{subsec:10.osc.family}

Let $W\in C_c^\infty(\R)$ be even. For parameters $(\tau,\Delta)\in\R\times(0,\infty)$ define
\begin{equation}\label{eq:10.htaudelta}
h_{\tau,\Delta}(t):=W(t/\Delta)\,e^{-i\tau t}.
\end{equation}
With the Fourier convention of Chapter~9,
\begin{equation}\label{eq:10.hhat}
\widehat{h_{\tau,\Delta}}(\xi)=\Delta\,\widehat W\bigl(\Delta(\xi-\tau)\bigr),
\end{equation}
so $h_{\tau,\Delta}$ localizes near frequency $\xi=\tau$ with resolution $\asymp \Delta^{-1}$.

\begin{definition}[Smoothed oscillatory trace observable]\label{def:10.Ddef}
Define
\begin{equation}\label{eq:10.Ddef}
\mathfrak D_\Gamma(\tau;\Delta):=\mathcal T_\Gamma[h_{\tau,\Delta}].
\end{equation}
\end{definition}

\subsection{Stability under admissible auxiliary changes}
\label{subsec:10.osc.stability}

\begin{proposition}[Stability of $\mathfrak D_\Gamma(\tau;\Delta)$]\label{prop:10.D.stability}
Let $\mathfrak D_\Gamma(\tau;\Delta)$ be defined by \eqref{eq:10.Ddef} using a fixed admissible construction, and let
$\mathfrak D'_\Gamma(\tau;\Delta)$ be defined using any other admissible construction (different cutoff, quantization,
microlocal partition, or regularization in the noncompact case). Then
\begin{equation}\label{eq:10.D.stability}
\mathfrak D_\Gamma(\tau;\Delta)-\mathfrak D'_\Gamma(\tau;\Delta)
=
\Delta_{\mathrm{id}}(\tau,\Delta)
+
\Err_\Gamma(\tau,\Delta),
\end{equation}
where $\Delta_{\mathrm{id}}(\tau,\Delta)$ is an explicit identity-sector correction and
\begin{equation}\label{eq:10.D.stability.bound}
|\Err_\Gamma(\tau,\Delta)|
\le
\mathcal E_\Gamma[h_{\tau,\Delta}]+\mathcal E'_\Gamma[h_{\tau,\Delta}].
\end{equation}
\end{proposition}

\begin{proof}
This is an immediate application of Theorem~\ref{thm:08.normalization-stability} to the test $h_{\tau,\Delta}$.
\end{proof}

% ======================================================================
\section{Zeta-type encoders: Euler products and logarithmic derivatives}
\label{sec:10.zeta-encoders}
% ======================================================================

\subsection{Geometric encoding via logarithmic derivatives}
\label{subsec:10.euler}

In many trace formula settings the geometric side reorganizes into a Laplace transform of length/arithmetic data.
A prototype is the Selberg zeta function in rank-one hyperbolic settings.

Let $\{\gamma\}$ range over primitive hyperbolic conjugacy classes and let $\ell(\gamma)$ denote their lengths.
Assume there exists a zeta-type function $Z(s)$, holomorphic for $\Re(s)\gg1$, such that
\begin{equation}\label{eq:10.logder}
-\frac{Z'(s)}{Z(s)}
=
\sum_{\{\gamma\}}\sum_{m\ge1} a(\gamma^m)\,e^{-sm\ell(\gamma)},
\qquad \Re(s)\gg1,
\end{equation}
with amplitude factors $a(\gamma^m)$ matching the geometric weights of the trace formula.

\subsection{Spectral encoding}
\label{subsec:10.spectral-interp}

In compact rank-one cases, the nontrivial zeros of $Z(s)$ encode Laplace eigenvalues; in noncompact settings, scattering poles
enter through additional factors (e.g.\ determinants of scattering matrices).
In all cases the trace formula provides a distributional identity matching a spectral sum/integral to a geometric sum.

After normalization, this becomes a canonical equality of the form
\begin{equation}\label{eq:10.explicit.schematic}
\mathcal T_\Gamma[h]=\mathcal G_\Gamma[h]+\mathcal S_\Gamma[h]+\Err_\Gamma[h],
\end{equation}
where $\mathcal G_\Gamma[h]$ is built from $\widehat h(\ell(\gamma))$ and $\mathcal S_\Gamma[h]$ is built from spectral parameters
(eigenvalues and, when present, scattering data), and
\[
|\Err_\Gamma[h]|\le \mathcal E_\Gamma[h].
\]

% ======================================================================
\section{Explicit formulae in normalized (portable) form}
\label{sec:10.explicit}
% ======================================================================

\subsection{A canonical template}
\label{subsec:10.explicit.template}

\begin{theorem}[Normalized explicit formula]\label{thm:10.normalized-explicit}
Let $h$ be admissible. Then
\begin{equation}\label{eq:10.explicit}
\mathcal T_\Gamma[h]=\mathcal G_\Gamma[h]+\mathcal S_\Gamma[h]+\Err_\Gamma[h],
\end{equation}
where:
\begin{enumerate}
\item the geometric contribution admits the representation
\begin{equation}\label{eq:10.G.form}
\mathcal G_\Gamma[h]=\sum_{\{\gamma\}}\sum_{m\ge1} a(\gamma^m)\,\widehat h\!\bigl(m\ell(\gamma)\bigr)
\end{equation}
whenever the setting admits the Euler/log-derivative decomposition \eqref{eq:10.logder};
\item the spectral contribution may be written as a sum or integral over spectral parameters $\rho$ (eigenvalues and, when present,
scattering poles/weights) of the form
\begin{equation}\label{eq:10.S.form}
\mathcal S_\Gamma[h]=\sum_{\rho} b(\rho)\,\widehat h(\rho)
\quad\text{or}\quad
\mathcal S_\Gamma[h]=\int b(\rho)\,\widehat h(\rho)\,d\rho,
\end{equation}
with coefficients $b(\rho)$ determined by multiplicities/scattering weights in the locked regularization;
\item the error term satisfies
\begin{equation}\label{eq:10.Err.bound}
|\Err_\Gamma[h]|\le \mathcal E_\Gamma[h].
\end{equation}
\end{enumerate}
\end{theorem}

\begin{proof}
This is the normalized trace decomposition from Chapter~8 applied to the trace formula of Chapter~7, with the
Euler/log-derivative rewriting \eqref{eq:10.logder} when available. The bound \eqref{eq:10.Err.bound} is by definition of
the consolidated ledger.
\end{proof}

\subsection{Smoothed explicit formula near a frequency}
\label{subsec:10.explicit.smoothed}

Applying Theorem~\ref{thm:10.normalized-explicit} to $h_{\tau,\Delta}$ yields the localized explicit formula
\begin{equation}\label{eq:10.explicit.smoothed.eq}
\mathfrak D_\Gamma(\tau;\Delta)=\mathcal G_\Gamma[h_{\tau,\Delta}]+\mathcal S_\Gamma[h_{\tau,\Delta}]+\Err_\Gamma(\tau,\Delta),
\end{equation}
with
\[
|\Err_\Gamma(\tau,\Delta)|\le \mathcal E_\Gamma[h_{\tau,\Delta}].
\]
In particular, $\mathcal S_\Gamma[h_{\tau,\Delta}]$ acts as a smoothed spectral density probe near $\tau$, while
$\mathcal G_\Gamma[h_{\tau,\Delta}]$ is a smoothed length-spectrum/arithmetic observable controlled by $\widehat W$.

% ======================================================================
\section{Quadratic forms and positivity criteria (positivity gates)}
\label{sec:10.positivity}
% ======================================================================

\subsection{A canonical nonnegative test construction}
\label{subsec:10.pos.conv}

Let $h$ be admissible and define the involution
\[
h^\ast(t):=\overline{h(-t)}.
\]
Define the convolution square
\begin{equation}\label{eq:10.convolution}
(h\star h^\ast)(t):=\int_\R h(t-u)\,h^\ast(u)\,du.
\end{equation}
Then
\begin{equation}\label{eq:10.convolution.fourier}
\widehat{h\star h^\ast}(\xi)=|\widehat h(\xi)|^2\ge0.
\end{equation}

\begin{definition}[Positivity quadratic form]\label{def:10.Qdef}
Define
\begin{equation}\label{eq:10.Q}
\mathcal Q_\Gamma[h]:=\mathcal T_\Gamma[h\star h^\ast].
\end{equation}
\end{definition}

\subsection{Positivity modulo explicit correction and ledger remainder}
\label{subsec:10.pos.gate}

\begin{theorem}[Positivity gate]\label{thm:10.pos-gate}
Let $h$ be admissible. Then there exists a spectral quantity $\mathcal P_\Gamma[h]\ge0$ such that
\begin{equation}\label{eq:10.posgate}
\mathcal Q_\Gamma[h]
=
\mathcal P_\Gamma[h]
+
\mathcal G_\Gamma[h\star h^\ast]
+
\Err_\Gamma^{\mathrm{pos}}[h],
\end{equation}
where $\mathcal G_\Gamma[h\star h^\ast]$ is the geometric term evaluated at the convolution square and
\begin{equation}\label{eq:10.pos.err}
\big|\Err_\Gamma^{\mathrm{pos}}[h]\big|
\le
\mathcal E_\Gamma[h\star h^\ast].
\end{equation}
In particular, once $\mathcal G_\Gamma[h\star h^\ast]$ is explicitly bounded and the ledger term is negligible in a given regime,
$\mathcal Q_\Gamma[h]$ is nonnegative up to an explicit, computable correction.
\end{theorem}

\begin{proof}
Evaluating the spectral side of the (normalized) trace formula on $h\star h^\ast$ produces the nonnegative weight
$|\widehat h|^2$ by \eqref{eq:10.convolution.fourier}; this defines $\mathcal P_\Gamma[h]\ge0$. The remaining contributions are the
geometric term and the ledger-bounded error.
\end{proof}

\subsection{Consequences: localization of spectral parameters}
\label{subsec:10.pos.localization}

In settings where spectral parameters correspond to zeros/poles of a zeta-type function, inequalities derived from
\eqref{eq:10.posgate} provide the analytic mechanism underlying Li--Weil type criteria: families of tests $h$ whose
$\widehat h$ concentrate near prescribed spectral windows convert control of the geometric term into constraints on admissible
spectral locations. Concretely, \eqref{eq:10.posgate} implies
\begin{equation}\label{eq:10.pos.localization.ineq}
\mathcal P_\Gamma[h]
\ge
-\big|\mathcal G_\Gamma[h\star h^\ast]\big|
-\mathcal E_\Gamma[h\star h^\ast],
\end{equation}
so any regime where the right-hand side is small forces $\mathcal P_\Gamma[h]$ (hence the spectral mass in the window selected by
$|\widehat h|^2$) to be small as well.

% ======================================================================
\section{Determinant identities and trace expansions}
\label{sec:10.det-trace}
% ======================================================================

\subsection{Formal determinant expansion}
\label{subsec:10.det.formal}

Let $K$ be trace class on $H$. Then
\begin{equation}\label{eq:10.det.id}
\log\det(I-K)=-\sum_{n\ge1}\frac{1}{n}\Tr(K^n),
\end{equation}
convergent for $\|K\|<1$ and extendable by analytic continuation in standard settings.

\subsection{Trace-to-zeta mechanism (portable formulation)}
\label{subsec:10.det.mechanism}

Assume there exists an operator family $K(s)$ such that:
\begin{enumerate}
\item for $\Re(s)\gg1$, $K(s)$ is trace class and $\|K(s)\|<1$;
\item for each $n\ge1$, $\Tr(K(s)^n)$ admits a geometric expansion
\begin{equation}\label{eq:10.Kn.trace.geom}
\Tr(K(s)^n)=\sum_{\{\gamma\}} c_n(\gamma)\,e^{-s\ell(\gamma)};
\end{equation}
\item the coefficients $\{c_n(\gamma)\}$ satisfy the multiplicativity needed to reorganize \eqref{eq:10.det.id} into a logarithmic
derivative/Euler product.
\end{enumerate}
Then \eqref{eq:10.det.id} yields
\begin{equation}\label{eq:10.euler.product}
\det(I-K(s))
=
\prod_{\{\gamma\}}
\exp\!\left(-\sum_{n\ge1}\frac{1}{n}c_n(\gamma)\,e^{-s\ell(\gamma)}\right),
\end{equation}
which in favorable cases coincides with a zeta-type function $Z(s)$.

\medskip
\noindent\textbf{Normalization clause.}
In the present work the admissible input traces in \eqref{eq:10.Kn.trace.geom} must be produced through stabilized trace objects
(Chapter~8): any auxiliary dependence is permitted only insofar as it is absorbed into (i) an explicit identity correction and
(ii) a quantified ledger remainder. This is the determinant-level analogue of portability for $\mathcal T_\Gamma$.

% ======================================================================
\section{Concluding remarks and portability checklist}
\label{sec:10.conclusion}
% ======================================================================

The chapter provides three interfaces transporting normalized trace information into zeta-type objects:

\begin{enumerate}
\item \textbf{Determinants:} zeta-regularized determinants are reconstructed from stabilized trace input
(\eqref{eq:10.mellin}--\eqref{eq:10.detA.def}) and determinant expansions \eqref{eq:10.det.id}, with admissible auxiliary dependence
confined to explicit identity corrections and ledger errors.
\item \textbf{Explicit formulae:} the normalized explicit formula \eqref{eq:10.explicit} packages the spectral--geometric dictionary
in a canonical form, with error bounded by $\mathcal E_\Gamma[h]$.
\item \textbf{Positivity:} convolution squares yield canonical quadratic forms \eqref{eq:10.Q} admitting a positivity decomposition
\eqref{eq:10.posgate}; this is the analytic mechanism behind Li--Weil type constraints.
\end{enumerate}

In the final chapter we assemble these bridges into the global pipeline: normalized trace $\Rightarrow$ explicit formula
$\Rightarrow$ determinant encoder $\Rightarrow$ positivity gates, and we list the remaining analytic tasks required to close
regime-by-regime ledger bounds in the specific geometric model under study.

% ======================================================================
\section*{Bibliographic notes}
% ======================================================================

For the Selberg trace formula and its connection to Selberg zeta, see \cite{Hejhal1983,Buser1992}.
For spectral zeta functions and zeta-regularized determinants, see \cite{Seeley1967,RaySinger1971}.
For semiclassical trace methods, resonances, and scattering determinants, see \cite{Zworski2012,DyatlovZworski2019}.
For positivity criteria related to zeta zeros (Li-type and related frameworks), see \cite{Li1997,Connes1999,Titchmarsh1986}.

% ======================================================================
% Audit (Chapter 10, Trace-to-zeta bridges)
% ======================================================================

\section*{Audit (Trace-to-zeta bridges)}\label{sec:10.audit}

\noindent\textbf{Locks.}
\begin{itemize}
\item[(L10.1)] Normalized trace object $\mathcal T_\Gamma$ fixed (Chapter~8); all changes $\Rightarrow$ explicit identity correction
$+\,$ledger remainder.
\item[(L10.2)] Bridge taxonomy: determinants / explicit formula / positivity gates are stated as objects with bounds.
\item[(L10.3)] Positivity: $\widehat{h\star h^\ast}=|\widehat h|^2\ge0$ is the only positivity input; everything else is explicit
geometric contribution + ledger.
\end{itemize}

\medskip
\noindent\textbf{Failure modes.}
\begin{itemize}
\item[(F10.1)] Hidden identity leakage (forgetting $h(0)\chi$ subtraction) $\Rightarrow$ non-portable determinant/positivity claims.
\item[(F10.2)] Unstated regularization of continuous spectrum $\Rightarrow$ object-swap between $\mathcal C$ and $\mathcal S$.
\item[(F10.3)] Using $\mathcal Q_\Gamma[h]\ge0$ without accounting for $\mathcal G_\Gamma[h\star h^\ast]$ and $\mathcal E_\Gamma$
$\Rightarrow$ false positivity.
\end{itemize}

\medskip
\noindent\textbf{STATUS\_CERT.}
CLOSED (structural): bridges stated in portable normalized form with explicit identity isolation and ledger-bounded errors.
NEXT(1): instantiate $\mathcal G_\Gamma$ and $\mathcal S_\Gamma$ for the specific geometric model (compact/noncompact) and record the
regime where $\mathcal E_\Gamma[h_{\tau,\Delta}]$ and $\mathcal E_\Gamma[h\star h^\ast]$ are negligible.

% ======================================================================
% End of Chapter 10
% ======================================================================
